% !TeX root = ../se320_structural_assurability.tex

% ==================================================
\section{Applying Structural Assurability to AI Architectures}
\label{aa.sec.applications}
% ==================================================

The preceding sections developed structural assurability as a claim-relative
property concerning the capacity of a system and its evidence-producing
environment to make assurance-relevant distinctions possible.
This section applies that framework to representative AI architectures.

The purpose is not to evaluate particular commercial systems or prescribe
a preferred architecture.
Instead, the examples examine whether the formal distinctions developed in
Sections~\ref{aa.sec.foundations} and~\ref{aa.sec.formal-model}
reveal meaningful differences among systems that may otherwise appear
functionally similar.

Three cases are considered.

The first demonstrates strict comparative dominance for a claim concerning
unauthorized consequential action.

The second demonstrates assurability equivalence for a claim that can be
resolved primarily through externally observable behavior.

The third demonstrates incomparability when different architectures enable
different assurance-relevant distinctions.

Together, the cases illustrate that structural assurability is claim-relative
and need not induce a universal ranking of architectures.

\subsection{Representative Architectures}
\label{aa.subsec.representative-architectures}

Consider three abstract AI system architectures.

They are intentionally simplified.
Each represents a class of evidence-producing structure rather than a
specific implementation.

\subsubsection{Opaque Service Architecture}

Let \(S_O\) denote an AI system exposed primarily through an input-output
interface.

The system accepts requests and returns outputs.
Standard application-level records may establish that a request occurred
and that a response was returned, but the assurance interface exposes little
or no trusted evidence concerning internal execution, tool use,
authorization decisions, retrieved information, intermediate state,
or causal provenance.

A simplified evidence surface is

\begin{equation}
  S_O:
  \qquad
  \{
  \text{input},
  \text{output},
  \text{request metadata}
  \}.
\end{equation}

This architecture is not assumed to be unsafe or unreliable.
It is simply structurally limited in the assurance-relevant distinctions
that can be made from its exposed evidence.

\subsubsection{Instrumented Agent Architecture}

Let \(S_A\) denote an agentic AI system with runtime instrumentation.

In addition to externally visible inputs and outputs, the system records
evidence concerning consequential execution events, such as

\begin{equation}
  \begin{aligned}
    S_A:\quad
    \{ &
    \text{input},
    \text{output},
    \text{tool invocations},
    \text{retrieval events}, \\
       &
    \text{authorization decisions},
    \text{memory updates},
    \text{state transitions},
    \text{provenance}
    \}.
  \end{aligned}
\end{equation}

The architecture therefore provides a larger evidence surface concerning
system execution.

However, the evidence may still be generated and preserved within
infrastructure controlled by the same operator or execution environment
as the agent itself.

\subsubsection{Assurability-Oriented Architecture}

Let \(S_R\) denote an architecture explicitly designed to preserve
assurance-relevant evidence.

It includes the runtime evidence available from \(S_A\) and adds
mechanisms such as

\begin{equation}
  \begin{aligned}
    S_R:\quad \{ &
    \text{instrumented execution evidence},
    \text{independent action observation},   \\
                 &
    \text{tamper-evident evidence preservation},
    \text{authenticated provenance},         \\
                 &
    \text{controlled replay or intervention},
    \text{explicit action-surface coverage}, \\
                 &
    \text{selected internal-state instrumentation}
    \}.
  \end{aligned}
\end{equation}

The term \emph{assurability-oriented} is descriptive rather than normative.
The architecture is designed to increase the range of assurance-relevant
distinctions that can be made under specified conditions.

It does not follow that the system satisfies any particular assurance claim.

\subsection{Case 1: Unauthorized Consequential Action}
\label{aa.subsec.case-unauthorized-action}

Consider the assurance claim

\begin{equation}
  \begin{aligned}
    C_1 ={} & \text{``The agent performed no unauthorized consequential action} \\
            & \text{during interval }T\text{.''}
  \end{aligned}
  \label{eq:claim-unauthorized-action}
\end{equation}

Suppose a consequential external action \(a\) is observed during \(T\).

Relevant assurance questions include:

\begin{align*}
  Q_1 & = \text{Did action }a\text{ occur?}                                                  \\
  Q_2 & = \text{Which component or actor initiated }a\text{?}                                \\
  Q_3 & = \text{What authorization state applied when }a\text{ was initiated?}               \\
  Q_4 & = \text{Through what execution pathway did }a\text{ occur?}                          \\
  Q_5 & = \text{Could the evidence concerning }a\text{ have been altered or suppressed?}     \\
  Q_6 & = \text{Is the evidence concerning }a\text{ corroborated by a source}                \\
      & \quad\text{outside the evaluated component?}                                         \\                                                        \\
  Q_7 & = \text{Were all consequential-action channels represented by the evidence surface?}
\end{align*}

These questions define claim-relevant distinctions.

For example, \(Q_2\) may require distinguishing

\begin{equation}
  \text{agent initiated }a
  \quad\text{from}\quad
  \text{human initiated }a,
\end{equation}

while \(Q_3\) may require distinguishing

\begin{equation}
  \text{authorized action}
  \quad\text{from}\quad
  \text{unauthorized action}.
\end{equation}

\subsubsection{Opaque Service}

Under \(S_O\), external evidence may establish that a result occurred,
but the evidence surface may not permit reliable attribution of the
action to the agent, a human operator, an external service, or another
component.

Thus, the architecture may permit resolution of \(Q_1\) while leaving
\(Q_2\)--\(Q_7\) substantially unresolved.

Its claim-relevant distinction set may therefore be represented
schematically as

\begin{equation}
  \mathcal{G}_{q_1}(S_O)
  =
  \{
  g_{\mathrm{occurrence}}
  \}.
\end{equation}

\subsubsection{Instrumented Agent}

System \(S_A\) additionally exposes runtime traces linking actions to
tool invocations, authorization decisions, retrieval events, and state
transitions.

It may therefore permit distinctions concerning:

\begin{equation}
  \begin{aligned}
     & \text{action occurrence},   \\
     & \text{action initiator},    \\
     & \text{authorization state}, \\
     & \text{execution pathway}.
  \end{aligned}
\end{equation}

Schematically,

\begin{equation}
  \mathcal{G}_{q_1}(S_A)
  =
  \{
  g_{\mathrm{occurrence}},
  g_{\mathrm{origin}},
  g_{\mathrm{authorization}},
  g_{\mathrm{pathway}}
  \}.
\end{equation}

However, if all records are generated and preserved within the same
administrative or computational boundary as the agent, questions
concerning evidence manipulation or independent corroboration may
remain unresolved.

\subsubsection{Assurability-Oriented Architecture}

System \(S_R\) preserves the distinctions available from \(S_A\) while
adding independently generated observations, tamper-evident records,
authenticated provenance, and explicit monitoring of consequential-action
channels.

It may therefore additionally distinguish:

\begin{equation}
  \text{unaltered evidence}
  \quad\text{from}\quad
  \text{evidence subject to undetected manipulation},
\end{equation}

and

\begin{equation}
  \text{complete monitored action surface}
  \quad\text{from}\quad
  \text{partially observed action surface}.
\end{equation}

Schematically,

\begin{equation}
  \begin{aligned}
    \mathcal{G}_{q_1}(S_R)
    =
    \{ &
    g_{\mathrm{occurrence}},
    g_{\mathrm{origin}},
    g_{\mathrm{authorization}},
    g_{\mathrm{pathway}}, \\
       &
    g_{\mathrm{integrity}},
    g_{\mathrm{coverage}},
    g_{\mathrm{independent\ corroboration}}
    \}.
  \end{aligned}
\end{equation}

If each architecture preserves the capabilities of the preceding
architecture while adding the capabilities shown above, then

\begin{equation}
  \mathcal{G}_{q_1}(S_O)
  \subset
  \mathcal{G}_{q_1}(S_A)
  \subset
  \mathcal{G}_{q_1}(S_R).
\end{equation}

Therefore,

\begin{equation}
  S_R
  \succ_{q_1}
  S_A
  \succ_{q_1}
  S_O.
  \label{eq:case1-order}
\end{equation}

This ordering is not a general ranking of the three architectures.

It holds only for the specified claim and assurance context.

The result illustrates how structural differences can create strict
comparative assurability without establishing that any of the systems
actually satisfies \(C_1\).

Indeed, \(S_R\) may provide stronger evidence that an unauthorized action
\emph{did} occur.

\subsection{Case 2: Externally Observable Output Property}
\label{aa.subsec.case-output-property}

Consider a different claim:

\begin{equation}
  \begin{aligned}
    C_2 ={} & \text{``The response returned for request }x\text{ contains no} \\
            & \text{instance of prohibited output class }P\text{.''}
  \end{aligned}
  \label{eq:claim-output}
\end{equation}

Assume that membership in \(P\) can be determined completely and
reliably from the externally returned response.

The material assurance question is then simply whether

\begin{equation}
  y \in P,
\end{equation}

where \(y\) is the externally observable output.

All three architectures expose \(y\).

Additional information concerning tool calls, internal execution,
provenance, independent runtime observation, or state transitions may
be valuable for other purposes, but under the assumptions of this claim
it enables no additional distinction required to determine whether
\(y\in P\).

Thus,

\begin{equation}
  \mathcal{G}_{q_2}(S_O)
  =
  \mathcal{G}_{q_2}(S_A)
  =
  \mathcal{G}_{q_2}(S_R).
\end{equation}

Therefore,

\begin{equation}
  S_O
  \sim_{q_2}
  S_A
  \sim_{q_2}
  S_R.
  \label{eq:case2-equivalence}
\end{equation}

This case is important because it prevents a common misinterpretation
of structural assurability.

More instrumentation does not necessarily imply greater assurability.

If the assurance-relevant distinction can already be resolved completely
from externally available evidence, additional structural evidence may
contribute nothing to that particular assurance context.

Thus,

\begin{equation}
  S_R\succ_{q_1}S_O
\end{equation}

and

\begin{equation}
  S_R\sim_{q_2}S_O
\end{equation}

can both hold without contradiction.

Assurability is claim-relative.

\subsection{Case 3: Internal Visibility Versus Independent Evidence}
\label{aa.subsec.case-incomparability}

The previous examples produced strict dominance and equivalence.

A third case illustrates incomparability.

Consider two architectures, \(S_V\) and \(S_N\).

System \(S_V\) provides extensive visibility into internal execution.
It exposes selected intermediate representations, internal state transitions,
detailed tool-selection processes, and model-generated explanations of
relevant actions.

However, these observations are produced and preserved by infrastructure
controlled by the same operator and execution environment as the system
under evaluation.

System \(S_N\), by contrast, exposes little internal state.
It provides only limited visibility into model computation.

However, consequential external actions are independently observed
by a separate enforcement or monitoring component.
Those records are authenticated and protected from retrospective
alteration by the evaluated system.

Consider the claim

\begin{equation}
  C_3 =
  \text{``The agent intentionally attempted an unauthorized consequential action.''}
  \label{eq:claim-intention-action}
\end{equation}

The claim contains at least two materially different components.

One concerns internal process:

\begin{equation}
  \begin{aligned}
    q_{\mathrm{internal}} ={} & \text{Did the agent enter a state or process}     \\
                              & \text{relevant to attempted unauthorized action?}
  \end{aligned}
\end{equation}

The other concerns external execution and attribution:

\begin{equation}
  \begin{aligned}
    q_{\mathrm{external}} ={} & \text{Did a consequential unauthorized action occur,} \\
                              & \text{and can it be independently attributed?}
  \end{aligned}
\end{equation}

System \(S_V\) may enable distinctions over internal computational
behavior unavailable from \(S_N\):

\begin{equation}
  g_{\mathrm{internal}}
  \in
  \mathcal{G}_{q_3}(S_V)
  \setminus
  \mathcal{G}_{q_3}(S_N).
\end{equation}

System \(S_N\), however, may enable independently corroborated
distinctions concerning consequential external action that \(S_V\)
cannot provide:

\begin{equation}
  g_{\mathrm{independent}}
  \in
  \mathcal{G}_{q_3}(S_N)
  \setminus
  \mathcal{G}_{q_3}(S_V).
\end{equation}

Accordingly,

\begin{equation}
  \mathcal{G}_{q_3}(S_V)
  \nsubseteq
  \mathcal{G}_{q_3}(S_N)
\end{equation}

and

\begin{equation}
  \mathcal{G}_{q_3}(S_N)
  \nsubseteq
  \mathcal{G}_{q_3}(S_V).
\end{equation}

Therefore,

\begin{equation}
  S_V
  \parallel_{q_3}
  S_N.
  \label{eq:case3-incomparable}
\end{equation}

Neither system dominates.

A scalar assurance score would require assigning a relative value
to internal visibility and independent external evidence.

The structural framework does not impose such a tradeoff.

That decision belongs to the specific assurance process, threat model,
or decision authority.

\subsection{Recasting the Cases as Claim-Relevant Distinctions}
\label{aa.subsec.case-distinctions}

The three cases can be interpreted through a common principle.

Let \(\Omega_C\) represent alternatives relevant to assurance claim \(C\).

Evidence contributes to structural assurability when it permits
distinctions among alternatives in \(\Omega_C\) that would otherwise
remain unresolved.

For Case~1, relevant alternatives include:

\begin{equation}
  \begin{aligned}
    \Omega_{C_1}
    =
    \{ &
    \text{no consequential action},     \\
       &
    \text{authorized agent action},     \\
       &
    \text{unauthorized agent action},   \\
       &
    \text{human-initiated action},      \\
       &
    \text{externally initiated action}, \\
       &
    \text{action with unreliable historical record}
    \}.
  \end{aligned}
\end{equation}

The richer architectures expose evidence permitting progressively more
of these alternatives to be distinguished.

For Case~2,

\begin{equation}
  \Omega_{C_2}
  =
  \{
  y\in P,
  y\notin P
  \}.
\end{equation}

The externally visible response already permits the relevant distinction.
Additional structural instrumentation need not increase assurability.

For Case~3, the relevant alternatives span both internal-process
hypotheses and externally observable action hypotheses.
Different architectures distinguish different subsets, producing
incomparability.

These examples motivate a more general interpretation:

\begin{quote}
  Structural assurability concerns the capacity of a system and
  its evidence-producing environment to make distinctions among
  alternatives material to an assurance claim.
\end{quote}

Under this interpretation, the structural dimensions introduced in
Section~\ref{aa.subsec.candidate-dimensions} are important because
removing one can make an assurance-relevant distinction unavailable.

\subsection{Relationship to the Structural Dimensions}
\label{aa.subsec.cases-dimensions}

The examples also illustrate how the candidate structural dimensions
affect assurance-relevant distinctions.

\paragraph{Observability.}
Observability determines whether relevant states, events, or effects
enter the available evidence surface at all.

\paragraph{Controllability.}
Controllability determines whether an evaluator can create or reproduce
conditions required to distinguish claim-relevant alternatives.

\paragraph{Traceability.}
Traceability permits observed events to be distinguished according
to origin, causal pathway, or transformation history.

\paragraph{Reconstructability.}
Reconstructability preserves distinctions that may be required after
the original execution has ended.

\paragraph{Independence.}
Independence permits distinctions between system self-report and
corroboration generated outside the component or actor being evaluated.

\paragraph{Integrity.}
Integrity permits distinctions between trustworthy preservation of
historical evidence and histories subject to undetected alteration
or suppression.

\paragraph{Evidence Coverage.}
Evidence coverage determines which portions of the claim-relevant state,
event, or action surface admit evidentiary distinctions at all.

The dimensions therefore need not be interpreted as components of
a numerical assurability score.

Their more fundamental role is to govern which claim-relevant
distinctions the evidence-producing architecture makes possible.

\subsection{Cost and Practical Assurability}
\label{aa.subsec.case-cost}

The examples above can also be extended to include resource constraints.

Suppose \(S_A\) and \(S_R\) both permit reconstruction of a consequential
execution.

In \(S_A\), reconstruction requires several days of specialist analysis
across heterogeneous logs.

In \(S_R\), authenticated provenance and preserved execution state permit
the same claim-relevant distinction to be established automatically
within minutes.

Under an unconstrained budget, both systems may enable the same
evidentiary capability.

Under a resource constraint \(B^\ast\),

\begin{equation}
  \operatorname{cost}
  \left(
  g_{\mathrm{reconstruction}},
  S_A,
  K
  \right)
  >
  B^\ast
\end{equation}

while

\begin{equation}
  \operatorname{cost}
  \left(
  g_{\mathrm{reconstruction}},
  S_R,
  K
  \right)
  \le
  B^\ast.
\end{equation}

Then

\begin{equation}
  g_{\mathrm{reconstruction}}
  \notin
  \mathcal{G}_{(C,K,B^\ast)}(S_A)
\end{equation}

while

\begin{equation}
  g_{\mathrm{reconstruction}}
  \in
  \mathcal{G}_{(C,K,B^\ast)}(S_R).
\end{equation}

The systems therefore differ in practical assurability despite
admitting similar evidence in principle.

This distinction is particularly important for deployed AI systems.
Assurance mechanisms that are technically possible but prohibitively
expensive, slow, fragile, or operationally disruptive may provide
little practical assurance under realistic constraints.

Structural features that make useful evidence inexpensive and routine
can therefore expand practical assurability without changing the
underlying assurance claim.

\subsection{What the Cases Demonstrate}
\label{aa.subsec.case-summary}

The three cases establish several properties of the framework.

First, systems exhibiting similar functional behavior can differ
materially in the assurance-relevant distinctions supported by their
evidence-producing structures.

Second, additional instrumentation does not automatically increase
assurability.
Its value depends upon the assurance claim.

Third, structural assurability is not a universal ranking of architectures.
The same pair of systems may exhibit strict dominance under one claim
and equivalence under another.

Fourth, architectures may be incomparable when each exposes
claim-material distinctions unavailable from the other.

Fifth, independence, integrity, provenance, internal visibility, and
coverage are not interchangeable forms of ``more evidence.''
They can enable qualitatively different assurance-relevant distinctions.

Sixth, resource constraints can convert differences in the efficiency
of evidence acquisition into differences in practical assurability.

Most importantly, the cases support the central interpretation
developed in this work:

\begin{quote}
  A system is structurally more assurable in a specified assurance
  context when its architecture and evidence-producing environment
  permit assurance-relevant distinctions that cannot practically
  be made for the comparison system, while preserving the
  distinctions already available from that system.
\end{quote}

This statement remains deliberately neutral regarding the eventual
assurance outcome.

A richer evidence structure may make it easier to support an
assurance claim, easier to refute it, or both.

Structural assurability concerns the capacity for examination,
not the conclusion of that examination.

The examples therefore reinforce the separation

\begin{equation}
  \boxed{
    \begin{aligned}
      \text{System Structure}
       & \rightarrow
      \text{Obtainable Evidence} \\
       & \rightarrow
      \text{Assurance-Relevant Distinctions}
      \rightarrow
      \text{Assurance Reasoning}.
    \end{aligned}
  }
\end{equation}

The present framework addresses the first three terms.

The final inference from evidence and distinctions to an assurance
conclusion remains the responsibility of the applicable assurance methodology.
